% \section{Síntese da revisão bibliográfica}
% \label{sec_sintese_rev_bib}

% A revisão bibliográfica evidencia que missões de \gls{rvdb} envolvem a interação de diferentes subsistemas, incluindo a dinâmica orbital relativa, o controle de atitude da espaçonave e a dinâmica de manipuladores robóticos utilizados em operações de captura.
% A literatura apresenta avanços relevantes em cada um desses tópicos, tanto na modelagem da dinâmica orbital quanto no estudo da dinâmica acoplada de sistemas robóticos espaciais e em estratégias de navegação e controle para operações de proximidade.

% Entretanto, observa-se que grande parte dos trabalhos existentes aborda esses aspectos de forma separada, analisando isoladamente a dinâmica orbital, o controle de atitude ou a dinâmica de manipuladores robóticos.
% Essa abordagem limita a representação das interações dinâmicas que ocorrem quando esses subsistemas operam simultaneamente durante missões de \gls{rvdb}.

% Nesse contexto, torna-se relevante o desenvolvimento de modelos capazes de integrar essas diferentes dinâmicas em um único ambiente de análise.
% Assim, esta dissertação propõe um modelo dinâmico que considera simultaneamente a dinâmica orbital relativa, a dinâmica acoplada entre a espaçonave e o manipulador robótico e o controle de atitude do sistema, permitindo investigar o comportamento do conjunto durante operações de aproximação e alinhamento em missões de \gls{rvdb}.